\begin{abstract}
  基准并非仅仅是一组测试题目的集合；它自然地在难度上诱导出一个经验测度。本文建立这种基于测度的基准评测观点。这一视角使我们能够识别并形式化两种核心失效模式：\emph{测度失配}，即基准的难度分布不同于预期的难度分布；以及\emph{覆盖缺口}，即难度轴上未被覆盖的区域使基准无法检测有意义的局部训练进展。我们提出两个诊断指标来刻画这两种失效模式，并给出严格的理论保证。在实验上，我们利用仔细估计的题目难度，分析了涵盖 \textbf{7} 个领域的 \textbf{25} 个常用基准。结果表明，许多基准隐含地过度加权某些难度区域，因此其排行榜结论应谨慎解读；另一些基准则在难度轴上留下大段未覆盖区间，从而限制了它们监测模型改进过程的能力。最后，我们将诊断转化为行动，提出 \textsc{BenchFill}，一种面向目标难度的简单题目改写算法，能够有效缩小局部覆盖缺口。
\end{abstract}
